\documentclass[11pt,oneside]{book}

% Cheat-sheet compile toggle.
% Uncomment the next line for compact cheat-sheet styling.
% \notescheatsheettrue
\newif\ifnotescheatsheet
\notescheatsheetfalse
\providecommand{\notescheatsheettoggle}{0}
\ifnum\notescheatsheettoggle=1\relax
  \notescheatsheettrue
\fi

% Page geometry: narrow left margin, generous right margin for notes.
% Cheatsheet mode uses tight uniform margins with no margin column.
\ifnotescheatsheet
\usepackage[
    paper=letterpaper,
    includeheadfoot,
    left=0.5in,
    right=0.5in,
    top=0.5in,
    bottom=0.5in,
    marginparsep=0pt,
    marginparwidth=0pt,
    headheight=15pt,
    headsep=0.12in,
    footskip=0.28in
]{geometry}
\else
\usepackage[
    paper=letterpaper,
    includeheadfoot,
    left=0.65in,
    right=2.58in,
    top=0.70in,
    bottom=0.85in,
    marginparsep=0.18in,
    marginparwidth=1.96in,
    headheight=15pt,
    headsep=0.18in,
    footskip=0.38in
]{geometry}
\fi

% Typography and microtypography.
\usepackage[T1]{fontenc}
\usepackage[utf8]{inputenc}
\usepackage{lmodern}
\usepackage{microtype}
\usepackage{textcomp}
\usepackage{setspace}
\ifnotescheatsheet
\setstretch{1.0}
\else
\setstretch{1.06}
\fi
\raggedbottom

% Core math and theorem support.
\usepackage{amsmath, amssymb, amsthm, mathtools}
\usepackage{dsfont}
\numberwithin{equation}{chapter}
\usepackage{mathrsfs}
\DeclareFontFamily{U}{rsfs}{\skewchar\font127 }
\DeclareFontShape{U}{rsfs}{m}{n}{<-6> rsfs5 <6-8> rsfs7 <8-> rsfs10}{}

% Tables, figures, colors, and graphics.
\usepackage{graphicx}
\usepackage{booktabs}
\usepackage{threeparttable}
\usepackage{array}
\usepackage{tabularx}
\usepackage{caption}
\usepackage{xcolor}
\usepackage{tikz}
\usetikzlibrary{calc}
\usepackage{pgfplots}
\pgfplotsset{compat=1.18}
\usepackage{soul}

% Boxes and marginalia.
\usepackage[most]{tcolorbox}
\tcbuselibrary{theorems}
\ifnotescheatsheet
\tcbset{
  lecrcodeboxstyle/.style={
    colback=white,
    colframe=black,
    boxrule=0.45pt,
    arc=0pt,
    left=4pt,
    right=4pt,
    top=4pt,
    bottom=4pt
  },
  lecnoteboxstyle/.style={
    arc=0pt,
    colback=white,
    colframe=black,
    left=3pt,
    right=3pt,
    top=3pt,
    bottom=3pt
  },
  lecthmbasetitlestyle/.style={
    arc=0pt,
    left=4pt,
    right=4pt,
    top=3pt,
    bottom=3pt
  },
  lecthmbasestyle/.style={
    arc=0pt,
    left=5pt,
    right=5pt,
    top=5pt,
    bottom=5pt
  }
}
\newcommand{\lecnotefont}{\tiny}
\else
\tcbset{
  lecrcodeboxstyle/.style={
    colback=gray!8,
    colframe=black!65,
    boxrule=0.45pt,
    arc=2.4mm,
    left=6pt,
    right=6pt,
    top=6pt,
    bottom=6pt
  },
  lecnoteboxstyle/.style={
    arc=2.4mm,
    colback=gray!4,
    colframe=black!65,
    left=5pt,
    right=5pt,
    top=4pt,
    bottom=4pt
  },
  lecthmbasetitlestyle/.style={
    arc=2.4mm,
    left=6pt,
    right=6pt,
    top=4pt,
    bottom=4pt
  },
  lecthmbasestyle/.style={
    arc=2.4mm,
    left=8pt,
    right=8pt,
    top=8pt,
    bottom=8pt
  }
}
\newcommand{\lecnotefont}{\footnotesize}
\fi
\usepackage{varwidth}
\usepackage{marginnote}
\usepackage{marginfix}
\setlength{\marginparpush}{18pt}
\renewcommand*{\marginfont}{\footnotesize}

% Links.
\usepackage{hyperref}
\hypersetup{
    colorlinks=true,
    linkcolor=black,
    citecolor=blue!70!black,
    urlcolor=blue!70!black,
    pdfauthor={Marcus Kuo},
    pdftitle={Notes}
}

% Lists and TOC helpers.
\usepackage{enumitem}
\usepackage{titling}
\usepackage{titlesec}
\usepackage{tocloft}
\usepackage{etoc}
\usepackage{xparse}
\renewcommand{\cfttoctitlefont}{\fontsize{24}{28}\selectfont\bfseries}\renewcommand{\cftaftertoctitle}{\hfill}\setlength{\cftbeforetoctitleskip}{0pt}\setlength{\cftaftertoctitleskip}{1.5em}

% Header/footer styling.
\usepackage{fancyhdr}
\usepackage{comment}
\newcommand{\notesauthor}{Marcus Kuo}
\newcommand{\notesinstitution}{The University of Chicago}
\newcommand{\notesyear}{2026}

\makeatletter
\renewcommand{\chaptermark}[1]{\markboth{Section \thechapter: #1}{}}
\makeatother

\fancypagestyle{notesbody}{%
    \fancyhf{}
    \fancyhead[L]{\itshape\nouppercase{\leftmark}}
    \fancyhead[R]{\itshape\notesauthor}
    \fancyfoot[L]{\itshape\notesyear}
    \fancyfoot[C]{\thepage}
    \fancyfoot[R]{\itshape\notesinstitution}
    \renewcommand{\headrulewidth}{0pt}
    \renewcommand{\footrulewidth}{0pt}
}

\fancypagestyle{plain}{%
    \fancyhf{}
    \fancyfoot[C]{\thepage}
    \renewcommand{\headrulewidth}{0pt}
    \renewcommand{\footrulewidth}{0pt}
}

\pagestyle{notesbody}

% Section headings.
\titleformat{\section}
  {\normalfont\bfseries\large}
  {\thesection}
  {0.9em}
  {}
\titleformat{\subsection}
  {\normalfont\bfseries\large}
  {\thesubsection}
  {0.8em}
  {}
\titlespacing*{\section}{0pt}{1.4\baselineskip}{0.8\baselineskip}
\titlespacing*{\subsection}{0pt}{1.1\baselineskip}{0.55\baselineskip}

% TOC appearance.
\setcounter{tocdepth}{2}
\renewcommand{\contentsname}{Contents}

% Figure/table captions: centered, caption above, notes below by command.
\captionsetup{
    font=normalsize,
    labelfont=bf,
    textfont=normalfont,
    justification=centering,
    singlelinecheck=false,
    skip=8pt
}
\captionsetup[figure]{position=top}
\captionsetup[table]{position=top}

% Code block styling.
\usepackage{listings}
\lstdefinelanguage{Rnote}{
  language=R,
  morekeywords={mutate,summarize,filter,group_by,mean,sd,replicate,ggplot,labs,geom_density}
}
\ifnotescheatsheet
\lstdefinestyle{notescode}{
  language=Rnote,
  basicstyle=\ttfamily\small,
  keywordstyle=\color{black},
  commentstyle=\itshape\color{black},
  stringstyle=\color{black},
  showstringspaces=false,
  columns=fullflexible,
  keepspaces=true,
  aboveskip=0pt,
  belowskip=0pt,
  frame=none
}
\else
\lstdefinestyle{notescode}{
  language=Rnote,
  basicstyle=\ttfamily\small,
  keywordstyle=\color{blue!75!black},
  commentstyle=\itshape\color{teal!55!black},
  stringstyle=\color{red!70!black},
  showstringspaces=false,
  columns=fullflexible,
  keepspaces=true,
  aboveskip=0pt,
  belowskip=0pt,
  frame=none
}
\fi
\newtcblisting{rcodebox}{
  enhanced,
  breakable,
  lecrcodeboxstyle,
  listing only,
  listing options={style=notescode}
}

% Margin note box.
\NewDocumentCommand{\lecnotebox}{ +m }{%
  \begin{tcolorbox}[
    enhanced,
    boxrule=0.45pt,
    lecnoteboxstyle,
    width=\marginparwidth,
    before skip=0pt,
    after skip=0pt
  ]
  {\lecnotefont #1}
  \end{tcolorbox}%
}

% Boxed-note backends. Normal mode uses \marginpar so marginfix can prevent
% overlap and bottom-margin spillover. Notes with an explicit offset fall back
% to \marginnote as a manual escape hatch.
\NewDocumentCommand{\lecnormalboxednoteplace}{ O{0pt} +m }{%
  \ifdim#1=0pt\relax
    \marginpar{\lecnotebox{#2}}%
  \else
    \marginnote{\lecnotebox{#2}}[#1]%
  \fi
}
\NewDocumentCommand{\leccheatsheetboxednoteplace}{ O{0pt} +m }{%
  \footnote{#2}%
}
\let\lecboxednoteplace\lecnormalboxednoteplace
\ifnotescheatsheet
\let\lecboxednoteplace\leccheatsheetboxednoteplace
\fi

% Notes triggered in math mode or inside theorem boxes are queued until we hit
% a stable text boundary where the chosen backend can place them safely.
\makeatletter
\newcommand{\lecdeferrednotes}{}
\NewDocumentCommand{\lecdefernote}{ O{0pt} +m }{%
  \g@addto@macro\lecdeferrednotes{\lecboxednoteplace[#1]{#2}}%
}
\newcommand{\lecflushdeferrednotes}{%
  \lecdeferrednotes
  \global\let\lecdeferrednotes\@empty
}
\AtEndDocument{\lecflushdeferrednotes}
\pretocmd{\appendix}{\lecflushdeferrednotes}{}{}
\makeatother

\NewDocumentCommand{\boxednote}{ O{0pt} +m }{%
  \ifmmode
    \lecdefernote[#1]{#2}%
  \else\ifinner
    \lecdefernote[#1]{#2}%
  \else\iflecinbox
    \lecdefernote[#1]{#2}%
  \else
    \lecflushdeferrednotes
    \lecboxednoteplace[#1]{#2}%
  \fi\fi\fi
}

% Margin figure helper: caption (with number) above, image nearly full width.
% Usage: \marginimage[Caption][fig:label]{filename}
% - `filename` is relative to `images/`
% - Omit caption/label by leaving the optional args empty
\NewDocumentCommand{\marginimage}{ O{} o m }{%
  \boxednote{%
    \begingroup
    \centering
    \IfNoValueTF{#2}{%
      \if\relax\detokenize{#1}\relax
        \includegraphics[width=\linewidth]{images/#3}%
      \else
        \refstepcounter{figure}%
        {\footnotesize\bfseries \figurename~\thefigure.\ \itshape #1\par}%
        \vspace{4pt}%
        \includegraphics[width=\linewidth]{images/#3}%
      \fi
    }{%
      \refstepcounter{figure}%
      \label{#2}%
      \if\relax\detokenize{#1}\relax
        \includegraphics[width=\linewidth]{images/#3}%
      \else
        {\footnotesize\bfseries \figurename~\thefigure.\ \itshape #1\par}%
        \vspace{4pt}%
        \includegraphics[width=\linewidth]{images/#3}%
      \fi
    }%
    \endgroup
  }%
}
\ifnotescheatsheet
\RenewDocumentCommand{\marginimage}{ O{} o m }{}%
\fi

% Notes under tables/figures.
\newcommand{\fignote}[1]{\par\vspace{0.25em}{\footnotesize\textit{Notes:} #1\par}}
\newcommand{\tabnote}[1]{\par\vspace{0.25em}{\footnotesize\textit{Notes:} #1\par}}

% Line-break penalties: prevent mid-expression line breaks.
\relpenalty=10000
\binoppenalty=10000

% =====================================================================
% MATH SHORTHAND MACROS
% =====================================================================
\newcommand*\N{\mathbb{N}}
\newcommand*\R{\mathbb{R}}
\newcommand*\Q{\mathbb{Q}}
\newcommand*\Z{\mathbb{Z}}
\newcommand*\E{\mathbb{E}}
\let\oldP\P
\renewcommand*\P{\mathbb{P}}
\let\O\varnothing
\let\oldemptyset\emptyset
\let\emptyset\varnothing
\newcommand{\ind}{\perp\!\!\!\perp}

\DeclareMathOperator*{\argmax}{arg\,max}
\DeclareMathOperator*{\argmin}{arg\,min}
\DeclareMathOperator*{\plim}{plim}
\DeclareMathOperator{\Var}{Var}
\DeclareMathOperator{\Cov}{Cov}
\DeclareMathOperator{\Corr}{Corr}
\DeclareMathOperator{\Bias}{Bias}
\DeclareMathOperator{\MSE}{MSE}

% =====================================================================
% CROSS-REFERENCE MACROS
% =====================================================================
\newcommand{\thmref}[1]{Theorem~\ref{thm:#1}}
\newcommand{\lemref}[1]{Lemma~\ref{thm:#1}}
\newcommand{\propref}[1]{Proposition~\ref{thm:#1}}
\newcommand{\defref}[1]{Definition~\ref{thm:#1}}
\newcommand{\corref}[1]{Corollary~\ref{thm:#1}}
\newcommand{\exref}[1]{Example~\ref{thm:#1}}

% =====================================================================
% THEOREM ENVIRONMENTS
% =====================================================================

% Shared counter: chapter.section.number (e.g., 1.2.3)
\newcounter{thmcounter}[section]
\renewcommand{\thethmcounter}{\thesection.\arabic{thmcounter}}

% Colors for theorem boxes.
% User palette (hex):
%  - theorem/proposition: #ccd5ae
%  - lemma/corollary:     #e9edc9 (different now)
%  - assumption:          #d4a373
%  - definition/notation: #fefae0
%  - example:             #faedcd (different now)
\definecolor{lecthm}{HTML}{CCD5AE}
\definecolor{leclem}{HTML}{ECF1D3}
\definecolor{lecassump}{HTML}{DDB892}
\definecolor{lecdef}{HTML}{FEFAE0}
\definecolor{lecex}{HTML}{FFD9C7}

% Main box fill shade: lower = darker (closer to the title bar color),
% higher = lighter (more white mixed in).
\newcommand{\lecmainmix}{45}

% Darken title bars + outlines slightly by mixing with black.
% Higher = darker.
\newcommand{\lectitlebarmix}{95}
\newcommand{\lecframemix}{95}
\ifnotescheatsheet
\colorlet{lecthm}{white}
\colorlet{leclem}{white}
\colorlet{lecassump}{white}
\colorlet{lecdef}{white}
\colorlet{lecex}{white}
\renewcommand{\lecmainmix}{100}
\renewcommand{\lectitlebarmix}{100}
\renewcommand{\lecframemix}{0}
\fi

% Proof-end symbol mechanism.
\newcommand{\currentproofmark}{\ensuremath{\square}}
\newcommand{\setproofmark}[1]{\renewcommand{\currentproofmark}{#1}}

% Default proof label (auto-updated by theorem-like environments).
\newcommand{\lecproofname}{Proof}
\newcommand{\setlecproofname}[1]{%
  \begingroup
  \edef\lecprooftmp{#1}%
  \global\let\lecproofname\lecprooftmp
  \endgroup
}
\newcommand{\resetlecproofname}{\setlecproofname{Proof}}

% Base tcolorbox style for theorem environments.
\tcbset{
  thmbasetitle/.style={
    enhanced,
    boxrule=0.45pt,
    lecthmbasetitlestyle,
    coltitle=black,
    fonttitle=\bfseries,
  },
  thmbase/.style={
    enhanced,
    breakable,
    boxrule=0.45pt,
    lecthmbasestyle,
    fonttitle=\bfseries,
    coltitle=black,
    separator sign none,
    description delimiters parenthesis,
    terminator sign={.},
    % Use \linewidth so titles wrap correctly inside lists/minipages.
    varwidth boxed title=\dimexpr\linewidth-3em\relax,
    attach boxed title to top left={xshift=7pt,yshift*=-2mm},
    boxed title style={thmbasetitle},
  },
	  thmbluestyle/.style={
	    thmbase,
	    colback=lecthm!\lecmainmix!white,
	    colframe=lecthm!\lecframemix!black,
	    boxed title style={thmbasetitle, colback=lecthm!\lectitlebarmix!black, colframe=lecthm!\lecframemix!black},
	  },
	  thmbgstyle/.style={
	    thmbase,
	    colback=leclem!\lecmainmix!white,
	    colframe=leclem!\lecframemix!black,
	    boxed title style={thmbasetitle, colback=leclem!\lectitlebarmix!black, colframe=leclem!\lecframemix!black},
	  },
	  thmyellowstyle/.style={
	    thmbase,
	    colback=lecdef!\lecmainmix!white,
	    colframe=lecdef!\lecframemix!black,
	    boxed title style={thmbasetitle, colback=lecdef!\lectitlebarmix!black, colframe=lecdef!\lecframemix!black},
	  },
	  thmpinkstyle/.style={
	    thmbase,
	    colback=lecex!\lecmainmix!white,
	    colframe=lecex!\lecframemix!black,
	    boxed title style={thmbasetitle, colback=lecex!\lectitlebarmix!black, colframe=lecex!\lecframemix!black},
	  },
	  thmassumpstyle/.style={
	    thmbase,
	    colback=lecassump!\lecmainmix!white,
	    colframe=lecassump!\lecframemix!black,
	    boxed title style={thmbasetitle, colback=lecassump!\lectitlebarmix!black, colframe=lecassump!\lecframemix!black},
	  },
	}

% Boxed theorem environments via newtcbtheorem.
% Usage: \begin{theorem}{Name}{label} ... \end{theorem}
%   - Name appears in parentheses after "Theorem 1.2.3"
%   - Empty name: \begin{theorem}{}{label}
%   - Label is prefixed with "thm:" automatically

\newtcbtheorem[use counter=thmcounter, number freestyle={\noexpand\thesection.\noexpand\arabic{thmcounter}}]{lectheorem}{Theorem}{
  thmbluestyle,
  before upper={
    \setproofmark{\ensuremath{\square}}
    \setlecproofname{Proof of Theorem~\thethmcounter}
  },
}{thm}

\newtcbtheorem[use counter=thmcounter, number freestyle={\noexpand\thesection.\noexpand\arabic{thmcounter}}]{lecproposition}{Proposition}{
  thmbluestyle,
  before upper={
    \setproofmark{\ensuremath{\square}}
    \setlecproofname{Proof of Proposition~\thethmcounter}
  },
}{thm}

\newtcbtheorem[use counter=thmcounter, number freestyle={\noexpand\thesection.\noexpand\arabic{thmcounter}}]{leclemma}{Lemma}{
  thmbgstyle,
  before upper={
    \setproofmark{\ensuremath{\blacksquare}}
    \setlecproofname{Proof of Lemma~\thethmcounter}
  },
}{thm}

\newtcbtheorem[use counter=thmcounter, number freestyle={\noexpand\thesection.\noexpand\arabic{thmcounter}}]{leccorollary}{Corollary}{
  thmbgstyle,
  before upper={
    \setproofmark{\ensuremath{\circlearrowright}}
    \setlecproofname{Proof of Corollary~\thethmcounter}
  },
}{thm}

\newtcbtheorem[use counter=thmcounter, number freestyle={\noexpand\thesection.\noexpand\arabic{thmcounter}}]{lecdefinition}{Definition}{
  thmyellowstyle,
}{thm}

\newtcbtheorem[use counter=thmcounter, number freestyle={\noexpand\thesection.\noexpand\arabic{thmcounter}}]{lecexample}{Example}{
  thmpinkstyle,
  before upper={\setproofmark{\ensuremath{\triangle}} },
}{thm}

\newtcbtheorem[use counter=thmcounter, number freestyle={\noexpand\thesection.\noexpand\arabic{thmcounter}}]{lecassumption}{Assumption}{
  thmassumpstyle,
}{thm}

\newtcbtheorem[use counter=thmcounter, number freestyle={\noexpand\thesection.\noexpand\arabic{thmcounter}}]{lecnotation}{Notation}{
  thmyellowstyle,
}{thm}

% Wrapper environments: optional title, optional label.
%   - \begin{theorem}[Title] ... \end{theorem} (auto-label; use \label{...} only if you want to refer)
%   - \begin{theorem}[Title][mylabel] ... \end{theorem} (sets \label{thm:mylabel})
\newcounter{lecautolabel}
\newif\iflecinbox
\makeatletter
\newcommand{\lecsetthmcurrentlabel}{%
  \protected@edef\@currentlabel{\thethmcounter}%
}
\makeatother
\NewDocumentEnvironment{theorem}{ O{} o }{%
  \global\lecinboxtrue
  \IfNoValueTF{#2}{%
    \stepcounter{lecautolabel}%
    \begin{lectheorem}{#1}{auto-\arabic{lecautolabel}}%
    \lecsetthmcurrentlabel
  }{%
    \begin{lectheorem}{#1}{#2}%
    \lecsetthmcurrentlabel
  }%
}{\end{lectheorem}\global\lecinboxfalse\lecflushdeferrednotes}
\NewDocumentEnvironment{proposition}{ O{} o }{%
  \global\lecinboxtrue
  \IfNoValueTF{#2}{%
    \stepcounter{lecautolabel}%
    \begin{lecproposition}{#1}{auto-\arabic{lecautolabel}}%
    \lecsetthmcurrentlabel
  }{%
    \begin{lecproposition}{#1}{#2}%
    \lecsetthmcurrentlabel
  }%
}{\end{lecproposition}\global\lecinboxfalse\lecflushdeferrednotes}
\NewDocumentEnvironment{lemma}{ O{} o }{%
  \global\lecinboxtrue
  \IfNoValueTF{#2}{%
    \stepcounter{lecautolabel}%
    \begin{leclemma}{#1}{auto-\arabic{lecautolabel}}%
    \lecsetthmcurrentlabel
  }{%
    \begin{leclemma}{#1}{#2}%
    \lecsetthmcurrentlabel
  }%
}{\end{leclemma}\global\lecinboxfalse\lecflushdeferrednotes}
\NewDocumentEnvironment{corollary}{ O{} o }{%
  \global\lecinboxtrue
  \IfNoValueTF{#2}{%
    \stepcounter{lecautolabel}%
    \begin{leccorollary}{#1}{auto-\arabic{lecautolabel}}%
    \lecsetthmcurrentlabel
  }{%
    \begin{leccorollary}{#1}{#2}%
    \lecsetthmcurrentlabel
  }%
}{\end{leccorollary}\global\lecinboxfalse\lecflushdeferrednotes}
\NewDocumentEnvironment{definition}{ O{} o }{%
  \global\lecinboxtrue
  \IfNoValueTF{#2}{%
    \stepcounter{lecautolabel}%
    \begin{lecdefinition}{#1}{auto-\arabic{lecautolabel}}%
    \lecsetthmcurrentlabel
  }{%
    \begin{lecdefinition}{#1}{#2}%
    \lecsetthmcurrentlabel
  }%
}{\end{lecdefinition}\global\lecinboxfalse\lecflushdeferrednotes}
\NewDocumentEnvironment{example}{ O{} o }{%
  \global\lecinboxtrue
  \IfNoValueTF{#2}{%
    \stepcounter{lecautolabel}%
    \begin{lecexample}{#1}{auto-\arabic{lecautolabel}}%
    \lecsetthmcurrentlabel
  }{%
    \begin{lecexample}{#1}{#2}%
    \lecsetthmcurrentlabel
  }%
}{\end{lecexample}\global\lecinboxfalse\lecflushdeferrednotes}
\NewDocumentEnvironment{assumption}{ O{} o }{%
  \global\lecinboxtrue
  \IfNoValueTF{#2}{%
    \stepcounter{lecautolabel}%
    \begin{lecassumption}{#1}{auto-\arabic{lecautolabel}}%
    \lecsetthmcurrentlabel
  }{%
    \begin{lecassumption}{#1}{#2}%
    \lecsetthmcurrentlabel
  }%
}{\end{lecassumption}\global\lecinboxfalse\lecflushdeferrednotes}
\NewDocumentEnvironment{notation}{ O{} o }{%
  \global\lecinboxtrue
  \IfNoValueTF{#2}{%
    \stepcounter{lecautolabel}%
    \begin{lecnotation}{#1}{auto-\arabic{lecautolabel}}%
    \lecsetthmcurrentlabel
  }{%
    \begin{lecnotation}{#1}{#2}%
    \lecsetthmcurrentlabel
  }%
}{\end{lecnotation}\global\lecinboxfalse\lecflushdeferrednotes}

% Unboxed environments: remark and aside (amsthm, sharing thmcounter).
\theoremstyle{remark}
\newtheorem{remark}[thmcounter]{Remark}
\newtheorem{aside}[thmcounter]{Aside}
\newtheorem{conj}[thmcounter]{Conjecture}

% Override proof environment to use \currentproofmark as QED symbol.
\makeatletter
\renewenvironment{proof}[1][\lecproofname]{%
  \par\pushQED{\qed}%
  \normalfont \topsep6\p@\@plus6\p@\relax
  \trivlist
  \item[\hskip\labelsep\itshape #1\@addpunct{.}]\ignorespaces
}{%
  \popQED\endtrivlist\@endpefalse
  \resetlecproofname
}
\makeatother
\renewcommand{\qedsymbol}{\currentproofmark}

% =====================================================================
% TITLE PAGE MACRO
% =====================================================================
\newcommand{\makelectitle}[3]{%
  \begin{titlepage}
    \thispagestyle{empty}
    \begin{tikzpicture}[remember picture,overlay]
      % Vertical rule.
      \draw[line width=0.6pt]
        ([xshift=0.78in,yshift=-0.78in]current page.north west) --
        ([xshift=0.78in,yshift=0.92in]current page.south west);

      % Title block (positions tuned to match the reference layout).
      \node[
        anchor=north west,
        inner sep=0pt,
        % Use some of the outer margin for long titles, but keep within the page.
        text width=\dimexpr\textwidth+\marginparsep+\marginparwidth-0.60in\relax,
        align=left
      ] at ([xshift=1.18in,yshift=-3.55in]current page.north west) {%
        \setstretch{1.0}%
        \fontsize{34}{36}\selectfont\bfseries
        \spaceskip=0.35em plus 0.04em minus 0.015em\relax
        \xspaceskip=\spaceskip
        \hyphenchar\font=-1\relax
        \raggedright
        % Prefer non-hyphenated breaks on long titles.
        \pretolerance=10000\relax
        \tolerance=10000\relax
        \emergencystretch=2em\relax
        \hyphenpenalty=10000\relax
        \exhyphenpenalty=10000\relax
        #1%
      };

      \node[
        anchor=north west,
        inner sep=0pt,
        text width=\textwidth,
        align=left
      ] at ([xshift=1.18in,yshift=-5.05in]current page.north west) {%
        \fontsize{17}{20}\selectfont\itshape #2%
      };

      \node[
        anchor=north west,
        inner sep=0pt,
        text width=\textwidth,
        align=left
      ] at ([xshift=1.18in,yshift=-6.40in]current page.north west) {%
        \fontsize{14}{16}\selectfont\scshape #3%
      };
    \end{tikzpicture}
  \end{titlepage}
  \clearpage
}

% =====================================================================
% CHAPTER OPENER MACRO
% =====================================================================
\newlength{\chaptextwid}
\newcommand{\chapteropening}[1]{%
  \thispagestyle{plain}
  \vspace*{0.10in}
  \noindent
  % Span across the main text column only.
  \setlength{\chaptextwid}{\textwidth}%
  \begin{tikzpicture}[x=1in,y=1in,line cap=rect,line join=miter]
    \pgfmathsetmacro{\rr}{\chaptextwid/1in}
    \pgfmathsetlengthmacro{\lecchapthin}{0.6pt}
    \pgfmathsetlengthmacro{\lecchaphalfthin}{0.5*\lecchapthin}
    \pgfmathsetlengthmacro{\lecchapbar}{2.4pt}
    % The y-values are the centerlines of the thin rules (use lengths to avoid rounding).
    \pgfmathsetlengthmacro{\lecchaptop}{0.18in}
    \pgfmathsetlengthmacro{\lecchapbot}{-0.44in}
    % Rule edges so bars meet the full rule thickness (no anti-aliased gaps at corners).
    \pgfmathsetlengthmacro{\lecchaptophi}{\lecchaptop+\lecchaphalfthin}
    \pgfmathsetlengthmacro{\lecchaptoplo}{\lecchaptop-\lecchaphalfthin}
    \pgfmathsetlengthmacro{\lecchapbothi}{\lecchapbot+\lecchaphalfthin}
    \pgfmathsetlengthmacro{\lecchapbotlo}{\lecchapbot-\lecchaphalfthin}
    \pgfmathsetlengthmacro{\lecchapleftpad}{0.06in}

    % Left header: "Chapter" and number.
    \node[anchor=west,font=\fontsize{26}{28}\selectfont] (chapword) at (\lecchapleftpad,\lecchaptop) {Section};
    \node[anchor=west,font=\fontsize{42}{42}\selectfont\bfseries] (chapnum) at ([xshift=0.10in]chapword.east) {\thechapter};

    % Right header: author.
    \node[anchor=east,font=\fontsize{14}{16}\selectfont\itshape] at (\rr,0.27) {\notesauthor};

    % Frame rules as filled rectangles so corners are pixel-perfect.
    \coordinate (lectopstart) at ([xshift=0.10in]chapnum.east);
    \path[fill=black] (lectopstart |- 0,\lecchaptoplo) rectangle (\rr,\lecchaptophi);
    \path[fill=black] (0.00,\lecchapbotlo) rectangle (\rr,\lecchapbothi);

    % Bars drawn last (on top) to avoid anti-aliased join gaps at the corners.
    \path[fill=black] (0.00,\lecchaptophi) rectangle (\lecchapbar,\lecchapbotlo);
    \path[fill=black] (\rr,\lecchaptophi) ++(-\lecchapbar,0) rectangle (\rr,\lecchapbotlo);

    % Chapter title: right-aligned on the bottom rule (do not cover the right bar).
    \node[
      anchor=base east,
      font=\fontsize{22}{24}\selectfont\bfseries,
      fill=white,
      inner xsep=0.03in,
      inner ysep=0.04in
    ] at (\rr-0.04,\lecchapbot) {%
      \begin{varwidth}{0.95\chaptextwid}\raggedleft #1\end{varwidth}%
    };
  \end{tikzpicture}
	  \vspace{0.28in}
	  {\parindent=0pt
	   \etocsettocstyle{}{}
	   \etocsetnexttocdepth{section}
   \etocsetstyle{section}
     {}
     {\makebox[2.2em][l]{\etocnumber}\etocname\nobreak\leaders\hbox to 0.7em{.\hss}\hfill\etocpage\par}
     {}
     {}
   \localtableofcontents
  }
  \vspace{0.35in}
}

% Convenience wrapper: manual chapter step + styled opener.
\newcommand{\styledchapter}[1]{%
  \lecflushdeferrednotes
  \cleardoublepage
  \refstepcounter{chapter}%
  \setcounter{section}{0}%
  \setcounter{subsection}{0}%
  \setcounter{equation}{0}%
  \markboth{Section \thechapter: #1}{}%
  \addcontentsline{toc}{chapter}{\protect\numberline{\thechapter}#1}%
  \chapteropening{#1}%
}
